\documentclass[11pt,a4paper]{article}
\usepackage[utf8]{inputenc}
\usepackage[T1]{fontenc}
\usepackage[english]{babel}
\usepackage{geometry}
\geometry{margin=2.5cm}
\usepackage{hyperref}
\usepackage{parskip}
\usepackage{booktabs}
\usepackage{enumitem}

\begin{document}

\begin{center}
  {\LARGE\bfseries WRO Future Engineers 2026}\\[0.4em]
  {\large Project Report --- Autonomous 4WD Smart Robot Car}\\[0.3em]
  \rule{\linewidth}{1pt}
\end{center}

\medskip
\noindent
\textbf{Team Name:} [INSERT TEAM NAME] \hfill \textbf{Country:} Morocco \\
\textbf{Members:} [Name 1], [Name 2], [Name 3] \hfill \textbf{Coach:} [INSERT COACH NAME]

\rule{\linewidth}{0.4pt}

\section*{1. Introduction}

Our team designed and built an autonomous robot car for the WRO Future Engineers 2026
challenge. The robot uses a Raspberry Pi 5 as its main controller, combined with a
camera, an ultrasonic sensor, and an IMU to navigate the track, avoid obstacles, and
complete the parking manoeuvre without any human input.

\section*{2. Hardware}

The robot is built on a 4WD aluminium chassis. The main components are:

\begin{itemize}[leftmargin=1.5em, itemsep=2pt]
  \item \textbf{Raspberry Pi 5} --- main controller, runs all software on-board.
  \item \textbf{Camera Module 3 Wide} --- 120° FOV, used for lane detection and pillar recognition.
  \item \textbf{HC-SR04 ultrasonic sensor} --- front-mounted, stops the robot if an obstacle is closer than 20 cm.
  \item \textbf{MPU-6050 IMU} --- measures heading and helps with precise turns.
  \item \textbf{TB6612FNG motor driver} --- controls the four DC motors independently.
  \item \textbf{3S LiPo battery} --- powers the motors; a 5 V buck converter supplies the Pi.
\end{itemize}

\section*{3. Software}

All code is written in Python. The program has three main parts:

\begin{enumerate}[leftmargin=1.5em, itemsep=2pt]
  \item \textbf{Vision} --- OpenCV detects track boundaries and red/green pillars using HSV colour segmentation.
  \item \textbf{Decision} --- combines camera output and IMU heading to compute a steering correction with a PD controller.
  \item \textbf{Motor control} --- converts the steering correction into PWM signals for the left and right motor pairs.
\end{enumerate}

Obstacle avoidance rule: red pillar $\rightarrow$ pass on the right; green pillar $\rightarrow$ pass on the left.
The robot detects pillars at around 80 cm to steer smoothly around them.

\section*{4. Results}

\begin{center}
\begin{tabular}{lcc}
  \toprule
  Scenario & Attempts & Successes \\
  \midrule
  Open round (3 laps, no obstacles) & 10 & [X] \\
  Obstacle round                    & 10 & [X] \\
  Parking manoeuvre                 & 10 & [X] \\
  \bottomrule
\end{tabular}
\end{center}

The robot completed laps consistently. The main issue observed was [brief description, e.g.\ occasional detection failure under direct sunlight].

\section*{5. Challenges}

\begin{itemize}[leftmargin=1.5em, itemsep=2pt]
  \item \textbf{Lighting changes} affected colour detection --- solved by adjusting HSV thresholds adaptively.
  \item \textbf{IMU drift} caused small heading errors --- solved by re-zeroing the yaw at each lap start line.
  \item \textbf{Motor asymmetry} made the robot drift sideways --- solved with a calibration offset stored in a config file.
\end{itemize}

\section*{6. Conclusion}

We built a working autonomous vehicle that navigates the WRO track, avoids coloured pillars, and completes the parking sequence. The project taught us computer vision, sensor integration, and real-time control. We plan to improve pillar detection with a small neural network for the next stage.

\end{document}